% ============================================================
%  Part II — Articles, possessives and case
% ============================================================
\gpart{III}{Articles \& Case}

% ------------------------------------------------------------
\gsec{The four cases at a glance}

\gintro{The case shows who (the person doing the action) does what to whom (the person receiving the action).}

\begin{gtbl}{colspec={X[0.85,l] X[1.35,l] X[0.75,l] X[1.5,l]}}
  Case & Job in the sentence & Asks & Example \\
  Nominativ & subject & wer? was? & \textbf{Der Hund} beißt den Mann. \\
  Akkusativ & direct object & wen? was? & Der Hund beißt \textbf{den Mann}. \\
  Dativ & indirect object & wem? & Ich gebe \textbf{dem Mann} das Buch. \\
  Genitiv & possessor & wessen? & das Auto \textbf{des Mannes} \\
\end{gtbl}

\tcap{The definite article across all four cases}

\begin{gendertbl}{}
  Definite & maskulin & feminin & neutrum & Plural \\
  Nominativ & der & die & das & die \\
  Akkusativ & \en{den} & die & das & die \\
  Dativ     & \en{dem} & \en{der} & \en{dem} & \en{den}\,+n \\
  Genitiv   & \en{des}\,+s & \en{der} & \en{des}\,+s & \en{der} \\
\end{gendertbl}

\begin{gnote}
The \textit{+n} and \textit{+s} marks are endings that go on the noun rather than the article.
\end{gnote}

% ------------------------------------------------------------
\gsec{Nominative}

\gintro{The nominative marks the subject — whoever or whatever is performing the
verb. It is also the form used after \textit{sein}, \textit{werden},
\textit{bleiben} and \textit{heißen}.}

\begin{posstbl}{}
  Group & Form & maskulin & feminin & neutrum & Plural \\
  \SetCell[r=3]{c,m} Articles & the   & der  & die   & das  & die   \\
                              & a     & ein  & ein\en{e}  & ein  & \xx   \\
                              & not a & kein & kein\en{e} & kein & kein\en{e} \\
  \hline\hline
  \SetCell[r=5]{c,m} Singular & my \eng{1st}      & mein   & mein\en{e}  & mein   & mein\en{e} \\
                              & your \eng{2nd}    & dein   & dein\en{e}  & dein   & dein\en{e} \\
                              & his \eng{3rd m}   & sein   & sein\en{e}  & sein   & sein\en{e} \\
                              & her \eng{3rd f}   & ihr    & ihr\en{e}   & ihr    & ihr\en{e} \\
                              & its \eng{3rd n}   & sein   & sein\en{e}  & sein   & sein\en{e} \\
  \SetCell[r=3]{c,m} Plural   & our \eng{1st}     & unser  & unser\en{e} & unser  & unser\en{e} \\
                              & your \eng{2nd}    & euer   & eur\en{e}   & euer   & eur\en{e} \\
                              & their \eng{3rd}   & ihr    & ihr\en{e}   & ihr    & ihr\en{e} \\
  Formal                      & your \eng{2nd}    & Ihr    & Ihr\en{e}   & Ihr    & Ihr\en{e} \\
\end{posstbl}

\begin{gnote}
The endings follow: \textit{-, -e, -, -e}.
\end{gnote}

\tcap{Examples}

\begin{extbl}{colspec={X[1,l] X[1.6,l]}}
  \textbf{Der} Hund schläft. & The dog is sleeping. \\
  \textbf{Meine} Schwester heißt Anna. & My sister is called Anna. \\
  Das ist \textbf{ein} gutes Buch. & That is a good book. \\
  \textbf{Kein} Mensch weiß das. & Nobody knows that. \\
\end{extbl}

% ------------------------------------------------------------
\gsec{Accusative}

\gintro{The accusative marks the direct object — the thing the verb acts on. It
also follows the accusative prepositions and measures spans of time.}

\begin{posstbl}{}
  Group & Form & maskulin & feminin & neutrum & Plural \\
  \SetCell[r=3]{c,m} Articles & the   & de\en{n}    & die   & das  & die   \\
                              & a     & eine\en{n}  & eine  & ein  & \xx   \\
                              & not a & keine\en{n} & keine & kein & keine \\
  \hline\hline
  \SetCell[r=5]{c,m} Singular & my \eng{1st}      & meine\en{n}  & meine        & mein         & meine \\
                              & your \eng{2nd}    & deine\en{n}  & deine        & dein         & deine \\
                              & his \eng{3rd m}   & seine\en{n}  & seine        & sein         & seine \\
                              & her \eng{3rd f}   & ihre\en{n}   & ihre         & ihr          & ihre \\
                              & its \eng{3rd n}   & seine\en{n}  & seine        & sein         & seine \\
  \SetCell[r=3]{c,m} Plural   & our \eng{1st}     & unsere\en{n} & unsere       & unser        & unsere \\
                              & your \eng{2nd}    & eure\en{n}   & eure         & euer         & eure \\
                              & their \eng{3rd}   & ihre\en{n}   & ihre         & ihr          & ihre \\
  Formal                      & your \eng{2nd}    & Ihre\en{n}   & Ihre         & Ihr          & Ihre \\
\end{posstbl}

\begin{gnote}
Only the masculine changes compared to the nominative.
\end{gnote}

\tcap{Examples}

\begin{extbl}{colspec={X[1,l] X[1.6,l]}}
  Ich sehe \textbf{den} Hund. & I see the dog. \\
  Sie kauft \textbf{einen} Tisch. & She is buying a table. \\
  Wir haben \textbf{keinen} Hunger. & We are not hungry. \\
  Er liest \textbf{seine} Zeitung. & He is reading his newspaper. \\
\end{extbl}

\begin{gnote}
The accusative also expresses duration and definite time without any
preposition: \textit{Ich bleibe \textbf{einen Monat}}, \textit{Er kommt
\textbf{jeden Tag}}, \textit{\textbf{Nächste Woche} fahre ich weg}.
\end{gnote}

% ------------------------------------------------------------
\gsec{Dative}

\gintro{The dative marks the indirect object — the person something is given to
or done for. It also follows the dative prepositions and a group of verbs whose
only object is dative.}

\begin{posstbl}{}
  Group & Form & maskulin & feminin & neutrum & Plural \\
  \SetCell[r=3]{c,m} Articles & the   & de\en{m}    & de\en{r}    & de\en{m}    & de\en{n}    \\
                              & a     & eine\en{m}  & eine\en{r}  & eine\en{m}  & \xx         \\
                              & not a & keine\en{m} & keine\en{r} & keine\en{m} & keine\en{n} \\
  \hline\hline
  \SetCell[r=5]{c,m} Singular & my \eng{1st}      & meine\en{m}  & meine\en{r}  & meine\en{m}  & meine\en{n} \\
                              & your \eng{2nd}    & deine\en{m}  & deine\en{r}  & deine\en{m}  & deine\en{n} \\
                              & his \eng{3rd m}   & seine\en{m}  & seine\en{r}  & seine\en{m}  & seine\en{n} \\
                              & her \eng{3rd f}   & ihre\en{m}   & ihre\en{r}   & ihre\en{m}   & ihre\en{n} \\
                              & its \eng{3rd n}   & seine\en{m}  & seine\en{r}  & seine\en{m}  & seine\en{n} \\
  \SetCell[r=3]{c,m} Plural   & our \eng{1st}     & unsere\en{m} & unsere\en{r} & unsere\en{m} & unsere\en{n} \\
                              & your \eng{2nd}    & eure\en{m}   & eure\en{r}   & eure\en{m}   & eure\en{n} \\
                              & their \eng{3rd}   & ihre\en{m}   & ihre\en{r}   & ihre\en{m}   & ihre\en{n} \\
  Formal                      & your \eng{2nd}    & Ihre\en{m}   & Ihre\en{r}   & Ihre\en{m}   & Ihre\en{n} \\
\end{posstbl}

\begin{gnote}
The endings follow \textit{-m, -r, -m, -n}.
\end{gnote}

\begin{gnote}
The dative plural also adds \textit{-n} to the noun:
\textit{die Kinder} $\rightarrow$ \textit{mit den Kinder\textbf{n}};
\textit{die Bücher} $\rightarrow$ \textit{in den Bücher\textbf{n}}. Plurals
already ending in \textit{-n} or \textit{-s} are left alone
(\textit{den Autos}).
\end{gnote}

\tcap{Examples}

\begin{extbl}{colspec={X[1,l] X[1.4,l]}}
  Ich gebe \textbf{dem} Kind einen Ball. & I give the child a ball. \\
  Sie hilft \textbf{ihrer} Mutter. & She helps her mother. \\
  Wir fahren mit \textbf{dem} Zug. & We are going by train. \\
  Das gehört \textbf{meinen} Eltern. & That belongs to my parents. \\
\end{extbl}

% ------------------------------------------------------------
\gsec{Genitive}

\gintro{The genitive marks possession, corresponding to English \emph{'s} or
\emph{of}. It also follows the genitive prepositions. In speech it is commonly
replaced by \textit{von + Dativ}, but it remains standard in writing.}

\begin{posstbl}{}
  Group & Form & maskulin & feminin & neutrum & Plural \\
  \SetCell[r=3]{c,m} Articles & the   & de\en{s}    & de\en{r}    & de\en{s}    & de\en{r}    \\
                              & a     & eine\en{s}  & eine\en{r}  & eine\en{s}  & \xx         \\
                              & not a & keine\en{s} & keine\en{r} & keine\en{s} & keine\en{r} \\
  \hline\hline
  \SetCell[r=5]{c,m} Singular & my \eng{1st}      & meine\en{s}  & meine\en{r}  & meine\en{s}  & meine\en{r} \\
                              & your \eng{2nd}    & deine\en{s}  & deine\en{r}  & deine\en{s}  & deine\en{r} \\
                              & his \eng{3rd m}   & seine\en{s}  & seine\en{r}  & seine\en{s}  & seine\en{r} \\
                              & her \eng{3rd f}   & ihre\en{s}   & ihre\en{r}   & ihre\en{s}   & ihre\en{r} \\
                              & its \eng{3rd n}   & seine\en{s}  & seine\en{r}  & seine\en{s}  & seine\en{r} \\
  \SetCell[r=3]{c,m} Plural   & our \eng{1st}     & unsere\en{s} & unsere\en{r} & unsere\en{s} & unsere\en{r} \\
                              & your \eng{2nd}    & eure\en{s}   & eure\en{r}   & eure\en{s}   & eure\en{r} \\
                              & their \eng{3rd}   & ihre\en{s}   & ihre\en{r}   & ihre\en{s}   & ihre\en{r} \\
  Formal                      & your \eng{2nd}    & Ihre\en{s}   & Ihre\en{r}   & Ihre\en{s}   & Ihre\en{r} \\
\end{posstbl}

\begin{gnote}
The endings run \textit{-s, -r, -s, -r}. Masculine and neuter match each other,
as do feminine and plural. The noun takes an ending too: \textit{-es} after one
syllable or after \textit{s, ß, z, tz, x} (\textit{des Kind\textbf{es}},
\textit{des Hau\textbf{ses}}), otherwise \textit{-s}
(\textit{des Lehrer\textbf{s}}).
\end{gnote}

\tcap{Examples}

\begin{extbl}{colspec={X[1,l] X[1.3,l]}}
  das Auto \textbf{des} Mannes & the man's car \\
  die Farbe \textbf{der} Tür & the colour of the door \\
  wegen \textbf{des} schlechten Wetters & because of the bad weather \\
  trotz \textbf{ihrer} Probleme & despite her problems \\
\end{extbl}

% ------------------------------------------------------------
\gsec{der-words and ein-words}

\gintro{Every determiner (the word in front of a noun) belongs to one of two families: der-words or ein-words. This determines its own ending and the adjective ending that follows it.}

\tcap{der-words, which take the definite-article endings}

\begin{flattbl}{colspec={X[0.85,l] X[1,l] X[0.85,l] X[1.1,l]}, vline{3}={1-Z}{0.5pt,solid,Line}, column{3}={font=\small\sffamily\bfseries, fg=Ink}}
  dieser  & this            & jeder    & each, every \eng{sg} \\
  jener   & that \eng{lit.} & mancher  & some, many a \\
  welcher & which           & solcher  & such \\
  aller   & all             & derselbe & the same \\
\end{flattbl}

\begin{gendertbl}{}
  dieser & maskulin & feminin & neutrum & Plural \\
  Nominativ & dies\en{er} & diese  & dies\en{es} & diese  \\
  Akkusativ & diesen & diese  & dies\en{es} & diese  \\
  Dativ     & diesem & dieser & diesem & diesen \\
  Genitiv   & dieses & dieser & dieses & dieser \\
\end{gendertbl}

\tcap{ein-words}

\begin{flattbl}{colspec={X[0.85,l] X[1,l] X[0.85,l] X[1.1,l]}, vline{3}={1-Z}{0.5pt,solid,Line}, column{3}={font=\small\sffamily\bfseries, fg=Ink}}
  ein  & a            & kein    & not a \\
  mein & my           & dein    & your \\
  sein & his, its     & ihr     & her \\
  unser & our         & euer    & your \\
  ihr & their         & Ihr     & your\\
\end{flattbl}

\begin{gendertbl}{}
  mein & maskulin & feminin & neutrum & Plural \\
  Nominativ & mein   & meine  & mein   & meine  \\
  Akkusativ & meinen & meine  & mein   & meine  \\
  Dativ     & meinem & meiner & meinem & meinen \\
  Genitiv   & meines & meiner & meines & meiner \\
\end{gendertbl}


